\section{Scaling the Pilot Paradox}
\label{sec:scaling}

The pilot refinement paradox is interesting only if it survives adequate
sampling. We therefore rerun the SQuAD/Mistral cell at \(n=500\) over five
seeds and three conditions: dense baseline, HCPC-v1 refinement, and HCPC-v2
gated refinement. Under the DeBERTa scorer, the baseline-minus-HCPC-v1 effect
is 1.1 faithfulness points, and only one of five seeds reaches \(p<0.05\) for
the paired contrast.

\begin{table}[H]
\centering
\small
\caption{Scaled SQuAD/Mistral headline cell. The pilot direction survives as a
small DeBERTa-measured non-monotonicity, not as a large stable effect.}
\label{tab:scaled}
\begin{tabular}{lrrrr}
\toprule
Condition & \(n\) & Faithfulness & Hallucination & Retrieval similarity \\
\midrule
Baseline & 2500 & 0.661 & 0.146 & 0.532 \\
HCPC-v1 & 2500 & 0.650 & 0.147 & 0.569 \\
HCPC-v2 & 2500 & 0.661 & 0.143 & 0.532 \\
\bottomrule
\end{tabular}
\end{table}

The result should be read as a scale correction. HCPC-v1 raises retrieval
similarity while lowering DeBERTa faithfulness, but the magnitude is small and
seed-sensitive. HCPC-v2 recovers the HCPC-v1 drop on average, yet this does
not establish a general solution to hallucination or a global method ranking.
